% fixture: amsthm/theorem-displays
% tier: T3
% engine: document
% covers: displays in theorems and proofs: the italic body's quad in \predisplaysize (which picks \abovedisplayshortskip or not, here at the margin), equation and equation* with amsmath, align in a proof, text resumed after a display in the body font; a proof ending with a display, whose \qed goes on a line of its own (the paragraph resumed after the display); a display mid-proof then text and the \qed
\documentclass{article}
\usepackage{amsmath,amsthm}
\newtheorem{theorem}{Theorem}
\begin{document}
\section{Results}
\begin{theorem}\label{thm:a}
For every integer $n$ we have
\begin{equation}\label{eq:sum}
\sum_{i=1}^n i = \frac{n(n+1)}{2},
\end{equation}
and the sum of squares is given by the next formula, which is a little longer:
\begin{equation*}
\sum_{i=1}^n i^2 = \frac{n(n+1)(2n+1)}{6}.
\end{equation*}
\end{theorem}
\begin{theorem}
One
\[ a\quad\quad\quad\quad\quad\quad\quad\quad\quad\quad\quad\quad\quad\;\;\;\:\!b \]
after.
\end{theorem}
\begin{proof}
By induction on $n$, using
\begin{align}
a &= b + c \\
  &= d
\end{align}
and \eqref{eq:sum}, the claim follows for all cases.
\end{proof}
\begin{proof}
It is enough to compute
\[ \sum_{i=1}^{n+1} i = \frac{n(n+1)}{2} + n + 1. \]
\end{proof}
Text after.
\end{document}
